% Zdroj: PDF priklady-leto-2025.pdf, fyzická str. 23-24. Značka: specializace UI-SU, UI-ZPJ. Zařazeno: S3.
\section{Shlukování}
\szzotazka{léto 2025}{31}{Shlukování (k-means)}

\noindent\textbf{Zadání.}\par
\begin{enumerate}
  \item Popište stručně algoritmus $k$-means. Jakým způsobem se volí počáteční polohy středů shluků (centroidů)? Jak se jejich
    polohy v průběhu algoritmu upravují? Kdy algoritmus končí?
  \item Jaké kritérium algoritmus $k$-means optimalizuje?
  \item Algoritmus $k$-means pro $k=3$ jsme spustili na data se třemi shluky z obrázku níže. Jednotlivé značky ($\times$, $\bullet$, $\blacksquare$) odlišují
    nalezené shluky. Proč algoritmus na těchto datech nenašel optimální shluky podle kritéria v předchozím bodě? Jak se
    dá zvýšit šance, že je algoritmus najde?
\end{enumerate}

\begin{center}
\includegraphics[width=0.7\linewidth]{shlukovani-leto2025.png}
\end{center}

\medskip
\noindent\textbf{Náčrt řešení.}\par
\begin{enumerate}
  \item Počáteční centroidy se mohou volit náhodně z dat. V průběhu algoritmu se každý bod z dat přiřadí k nejbližšímu
    centroidu, ten se potom přepočítá jako střed bodů k němu přiřazených. Tyto kroky se opakují dokud se mění přiřazení
    bodů ke shlukům, nebo po daný počet iterací.
  \item Minimalizuje se
    \[
      \sum_{i=1}^{N}\sum_{k=1}^{K} z_{i,k}\,\|x_i - \mu_k\|^2,
    \]
    kde $N$ je počet bodů, $K$ je počet středů a $z_{i,k}$ vyjadřuje, že bod $x_i$ patří do shluku se středem $\mu_k$ ($z_{i,k}=1$ pokud tam
    patří, jinak $0$).
  \item Prostřední shluk, který byl přiřazen k části levého, je mnohem menší než ostatní dva shluky. Šance na lepší výsledek
    se dá zlepšit opakováním náhodné inicializace počátečních centroidů a výběrem nejlepšího výsledku podle funkce v
    předchozím bodě.
\end{enumerate}

\medskip
\noindent\textit{Doplňující vysvětlení (nad rámec oficiálního náčrtu):} Na obrázku jsou tři přirozené shluky: velký dole vlevo, malý nahoře uprostřed a~jeden vpravo. Nalezené řešení rozdělilo velký levý shluk mezi dva středy ($\blacksquare$ a~$\bullet$) a~malý horní shluk připojilo ke středu $\bullet$. Takové rozdělení je \emph{lokálním} minimem kritéria: žádný bod nemá blíž k~jinému středu a~průměry se nemění, takže se algoritmus zastaví, ačkoli rozdělení podle přirozených shluků má menší součet čtverců. $k$-means konverguje jen k~lokálnímu minimu a~výsledek závisí na počáteční volbě středů. Kromě opakovaných běhů s~výběrem nejmenšího kritéria pomáhá inicializace \texttt{k-means++}, která rozprostře počáteční středy.
